\documentclass[twocolumn]{aastex631}
\begin{document}
\title{A residual reionization ionizing-photon-budget shortfall for star-forming galaxies at z\~{}9}
\author{NebulaMind Lab (autonomous pipeline)}
\affiliation{NebulaMind Lab, \url{https://lab.nebulamind.net}}
\begin{abstract}
Reionization ionizing-photon-budget reconciliation at z\~{}9: star-forming galaxies require f\_esc=0.390 (+0.393/-0.200) to close the budget (Madau-Dickinson SFRD, log xi\_ion=25.5+/-0.15, clumping C=2-5, JWST-SFRD tail), versus indirect-proxy-inferred f\_esc=0.062 (+0.108/-0.039) from LzLCS O32/beta calibrations. Median delta(required-inferred)=+0.302 dex-frac (16-84\%: +0.087 to +0.697); 93\% of the systematic MC shows a shortfall. Conclusion: a genuine SHORTFALL remains, and the sign holds under both O32 and beta calibrations. The result is bounded by the xi\_ion x clumping x proxy-calibration systematic, not by statistics. Generated autonomously from public data (jwst) via the NebulaMind Lab runner: a bounded, reproducible, descriptive study.
\end{abstract}
\section{Introduction}
Recent studies have highlighted a potential crisis in the photon budget during reionization [Muñoz2024]. This has led to increased demands on ionizing sources, questioning whether star-forming galaxies alone can account for the required ionizing photons [Davies2021]. To address this issue, we revisit the ionizing-photon-budget using a literature-anchored approach. Previous efforts have focused on excursion set reionization models and galaxy ionizing photon budgets at z < 10 [Park2022, Duncan2015], but our analysis aims to provide a more comprehensive understanding of the problem.
\section{Data and method}
Our method relies on published values for the cosmic star formation rate density (SFRD) from Madau \& Dickinson (2014), as well as calibrations for ionizing photon production efficiency (xi\_ion) and escape fraction (f\_esc) proxies. Specifically, we adopt xi\_ion = 10\^{}25.5 +/- 0.15 and O32/beta f\_esc proxy calibrations from LzLCS [Chisholm+22, Flury+22; Simmonds+24]. We do not utilize survey catalog data or observational data from JWST, SDSS, or TNG in this analysis.
\section{Result}
Our reconciliation of the reionization ionizing-photon-budget at z\~{}9 reveals that star-forming galaxies require an escape fraction f\_esc = 0.390 (+0.393/-0.200) to close the budget. This is in contrast to the indirect-proxy-inferred value of f\_esc = 0.062 (+0.108/-0.039) derived from LzLCS O32/beta calibrations. The median difference between the required and inferred escape fractions is +0.302 dex-frac, with a range of +0.087 to +0.697 (16-84\% confidence interval). Notably, 93\% of our systematic Monte Carlo simulations show a shortfall in ionizing photons.
\begin{figure}\centering\includegraphics[width=\columnwidth]{result.png}\caption{A residual reionization ionizing-photon-budget shortfall for star-forming galaxies at z\~{}9}\end{figure}
\section{Caveats}
It is essential to acknowledge the limitations of this analysis. Our results are based on an automated, single-selection, and uncalibrated measurement, which may introduce biases or inaccuracies. The reliance on published values for xi\_ion, clumping factor (C), and f\_esc proxy calibrations means that our findings are subject to the uncertainties inherent in these parameters. Additionally, our approach does not account for potential variations in galaxy properties or environmental factors that could impact the ionizing photon budget. Therefore, while our study provides valuable insights into the reionization crisis, further research is needed to refine and validate these results.
\section*{References}
\footnotesize
[Muoz2024] Muñoz et al. (2024). Reionization after JWST: a photon budget crisis?. bibcode 2024MNRAS.535L..37M \par [Davies2021] Davies et al. (2021). The Predicament of Absorption-dominated Reionization: Increased Demands on Ionizing Sources. bibcode 2021ApJ...918L..35D \par [Park2022] Park et al. (2022). Calibrating excursion set reionization models to approximately conserve ionizing photons. bibcode 2022MNRAS.517..192P \par [Duncan2015] Duncan et al. (2015). Powering reionization: assessing the galaxy ionizing photon budget at z < 10. bibcode 2015MNRAS.451.2030D \par [Madau2017] Madau (2017). Cosmic Reionization after Planck and before JWST: An Analytic Approach. bibcode 2017ApJ...851...50M

\end{document}
